%% Moving stuff here to clean up the main doc. Don't feel like digging around in past versions to find these agian if I want them.

%% DDrop Notes:

% unclear if "go" means "axe" or "should be written as"; parking it here
\hl{This can probably go: Based on our transcriptional profiles and on prior work showing that disruptions in IIS signaling can delay age-related decline in negative geotaxis}\cite{jones2009}, \hl{one might predict that prolonged R2 activation would improve climbing performance, as it induces downstream effectors of decreased insulin signaling} (\textit{Thor, foxo}).


SP: I tried to rephrase this paragraph, but I guess I didn't understand all of it:
Based on our lifespan and sequencing data, we hypothesized that green light may itself be a stressor, and that R4d>GtACR1 flies may endure this stress better than controls. If so, this mechanism did not manifest in female negative geotaxis data (\figref{fig:ddropfig}C). After one week of treatment, female R4d>GtACR1 flies covered less upward distance than controls (Light $\times$ Genotype p = 0.035), although we observed no difference in height climbed by 9 seconds. 
This difference appeared to be transient, as it was no longer detected after long-term inhibition. Males, however, outperformed controls after 1 week of inhibition (Light $\times$ Genotype p < 0.01 for height at 9 sec. and for total up distance). This advantage in both climbing metrics was retained after 5 weeks of inhibition (p < 0.01 for each). Effects of R4d activation were transient, by contrast (Supp XX). These data are consistent with R4d inhibition, but not activation, mitigating the impact of treatment on male climbing performance.